\documentclass[a4paper,11pt]{article}
\usepackage[T2A]{fontenc}
\usepackage[utf8]{inputenc}
\usepackage[russian]{babel}
\usepackage{amsmath,amssymb,mathtools}
\usepackage{icomma}
\usepackage{geometry}
\geometry{left=21mm,right=21mm,top=21mm,bottom=21mm}
\usepackage{microtype}
\usepackage{setspace}
\setstretch{1.12}
\usepackage{enumitem}
\usepackage{tabularx,booktabs,array}
\usepackage{tikz}
\usetikzlibrary{arrows.meta,positioning,shapes.geometric,calc}
\usepackage{xcolor}
\usepackage{fancyhdr}
\usepackage{hyperref}
\definecolor{accent}{HTML}{007E84}
\definecolor{darkaccent}{HTML}{15565B}
\definecolor{softgray}{HTML}{575E62}
\hypersetup{colorlinks=true,linkcolor=darkaccent}
\setlength{\parindent}{0pt}
\setlength{\parskip}{4.5pt}
\setlist[itemize]{leftmargin=1.5em,itemsep=2pt,topsep=3pt}
\setlist[enumerate]{leftmargin=1.8em,itemsep=5pt,topsep=3pt}
\newcommand{\key}[1]{\textcolor{darkaccent}{\textbf{#1}}}
\newcommand{\sepheading}[1]{\section{#1}\vspace{-4pt}\noindent\textcolor{accent}{\rule{\linewidth}{0.45pt}}\vspace{1pt}}
\renewcommand{\thesection}{\arabic{section}}
\makeatletter\@ifundefined{titleformatundefined}{}{}\makeatother
\pagestyle{fancy}
\fancyhf{}
\fancyhead[L]{\small\color{softgray}Финансовые задачи}
\fancyhead[R]{\small\color{softgray}Анастасия Павлова}
\fancyfoot[C]{\small\thepage}
\renewcommand{\headrulewidth}{0.4pt}
\renewcommand{\headrule}{\hbox to\headwidth{\color{accent}\leaders\hrule height \headrulewidth\hfill}}
\setlength{\headheight}{24pt}
\setlength{\headsep}{15pt}
\begin{document}
\begin{titlepage}
\thispagestyle{empty}
\begin{center}
\vspace*{31mm}
{\large\color{darkaccent}МАТЕМАТИКА}\par
\vspace{17mm}
{\LARGE Чек-лист}\par
\vspace{12mm}
{\LARGE\color{accent}Кредиты и вклады: схемы,\\[2mm] формулы, расчёты}\par
\vspace{15mm}
{\normalsize Подготовка к ЕГЭ по профильной математике}\par
\vfill
{\large Анастасия Павлова}\par
\vspace{12mm}
{\large 29 октября 2026 г.}\par
\vspace{20mm}
{\small Лёвин Артём Александрович}\par
{\small эксклюзивно для Анастасии Павловой}\par
\vspace{19mm}
\begin{tikzpicture}
\draw[accent,line width=0.9pt] (-5,0) .. controls (-3.4,0.75) and (-2.0,-0.6) .. (0,0)
 .. controls (1.8,0.63) and (3.4,-0.48) .. (5,0);
\fill[accent] (0,0) circle (1.15pt);
\end{tikzpicture}
\end{center}
\end{titlepage}

\sepheading{Как распознать схему кредитования}
В задачах о кредитах сначала устанавливаем, \key{какая величина остаётся одинаковой}.

\begin{tabularx}{\linewidth}{@{}p{0.28\linewidth} X@{}}
\toprule
\textbf{Схема} & \textbf{Узнаваемый признак в условии}\\
\midrule
Аннуитетная & Каждый период заёмщик вносит \textbf{равный платёж}. Тело долга при этом обычно убывает на разные суммы.\\[3pt]
Дифференцированная & Тело долга за каждый период уменьшается \textbf{на одну и ту же сумму}. Полный платёж с течением времени снижается.\\[3pt]
Произвольная & Задан иной порядок погашения, например долг ежегодно становится \textbf{на 25\% меньше предыдущего}. Размер погашаемой части в рублях меняется.\\
\bottomrule
\end{tabularx}

\key{Важно.} Слова «на одну и ту же сумму меньше» означают одинаковое \emph{уменьшение в рублях}. Слова «на 25\% меньше» означают уменьшение в одинаковое \emph{число раз} от текущего остатка. Эти условия различаются.

\sepheading{Дифференцированный кредит: пять формул}
Пусть $S$ --- исходная сумма кредита, $n$ --- число годовых периодов, $p=\frac{r}{100}$ --- годовая ставка в виде десятичной дроби. Проценты ежегодно начисляются \emph{на остаток долга}, а тело кредита погашается равными частями $S/n$.

\begin{align*}
\text{Общая сумма выплат:}\qquad
T&=S+Sp\,\frac{n+1}{2}.\\[2pt]
\text{Переплата:}\qquad
P&=Sp\,\frac{n+1}{2}.\\[2pt]
\text{Переплата в процентах от }S:\qquad
q&=100p\,\frac{n+1}{2}\,\%.\\[2pt]
\text{Первый платёж:}\qquad
A_1&=\frac Sn+Sp.\\[2pt]
\text{Последний платёж:}\qquad
A_n&=\frac Sn+\frac Sn\,p.
\end{align*}

\textbf{Почему в формуле появляется $\boldsymbol{(n+1)/2}$?} Сумма остатков, на которые начисляются проценты, равна
\[
S+\frac{n-1}{n}S+\cdots+\frac Sn
=\frac Sn(1+2+\dots+n)=S\frac{n+1}{2}.
\]
Слагаемое $Sp(n+1)/2$ --- суммарные проценты, начисленные за весь срок.

\textcolor{darkaccent}{\textbf{Пример с занятия.}} Кредит $1{,}5$ млн руб. на $4$ года под $15\%$ с равномерным уменьшением долга:
\[
T=1{,}5+1{,}5\cdot0{,}15\cdot\frac{4+1}{2}
=1{,}5+0{,}5625=2{,}0625\ \text{млн руб.}
\]
\key{Ответ:} $2\,062\,500$ руб.

\clearpage
\thispagestyle{fancy}
\begin{center}
{\Large\color{accent}Алгоритм: как выбрать схему кредита}
\end{center}
\vspace{17mm}
\begin{center}
\begin{tikzpicture}[
 node distance=12mm and 10mm,
 every node/.style={font=\normalsize,align=center},
 start/.style={ellipse,draw=darkaccent,very thin,minimum width=38mm,minimum height=11mm},
 decision/.style={diamond,aspect=2.6,draw=accent,very thin,text width=44mm,inner sep=3mm},
 action/.style={rounded corners=2mm,draw=darkaccent,very thin,text width=42mm,minimum height=15mm},
 arr/.style={-{Latex[length=2.6mm]},thick,color=darkaccent}
]
\node[start] (s) {Читаем условие};
\node[decision,below=of s] (eq) {Указаны равные\\ платежи?};
\node[action,right=14mm of eq] (ann) {Аннуитетная\\ схема};
\node[decision,below=25mm of eq] (dec) {Тело долга уменьшается\\ на равную сумму?};
\node[action,right=14mm of dec] (diff) {Дифференцированная\\ схема};
\node[action,below=22mm of dec] (arb) {Иной закон погашения:\\ работаем с условием\\ как с произвольной схемой};
\draw[arr] (s) -- (eq);
\draw[arr] (eq.east) -- node[above,font=\small]{Да} (ann.west);
\draw[arr] (eq) -- node[right,font=\small]{Нет} (dec);
\draw[arr] (dec.east) -- node[above,font=\small]{Да} (diff.west);
\draw[arr] (dec) -- node[right,font=\small]{Нет} (arb);
\end{tikzpicture}
\end{center}
\clearpage

\sepheading{Вклад с капитализацией и дополнительным взносом}
При капитализации начисленные проценты остаются на вкладе и тоже приносят проценты в следующий год.

Обозначения: $S$ --- первоначальная сумма; $A$ --- дополнительный взнос после первого года; $r$ --- ставка в процентах годовых; $k$ --- коэффициент роста:
\[
\boxed{\displaystyle k=1+\frac{r}{100}},\qquad k\ge1\ \text{при}\ r\ge0.
\]

\textbf{За два года} с пополнением после первого начисления процентов:
\[
\begin{aligned}
\text{после первого года:}&\quad Sk,\\
\text{после внесения суммы }A:&\quad Sk+A,\\
\text{после второго года:}&\quad (Sk+A)k=\boxed{Sk^2+Ak}.
\end{aligned}
\]

Если после первого года со вклада \textbf{снимают} $A$ рублей, получаем $Sk^2-Ak$ (при условии, что денег для снятия достаточно).

Если вклад длится три года, а после первого и второго годов вносится по $A$ рублей, итоговая сумма равна
\[
Sk^3+Ak^2+Ak.
\]
\key{Проверка смысла:} у первоначальной суммы $S$ две капитализации за два года, а у внесённой через год суммы $A$ --- только одна.

\sepheading{Как найти неизвестную процентную ставку}
Для двухлетнего вклада с пополнением составьте уравнение
\[
Sk^2+Ak=F,
\]
где $F$ --- известная итоговая сумма. Затем:
\begin{enumerate}
\item Перенесите всё в левую часть: $Sk^2+Ak-F=0$.
\item При необходимости избавьтесь от десятичных дробей, \emph{умножая каждый член} на $10$, $100$ и т.\,д.
\item Найдите дискриминант $D=A^2+4SF$ и корни по формуле.
\item Оставьте экономически допустимое значение $k\ge1$ (для неотрицательной ставки).
\item Переведите коэффициент в проценты: $r=100(k-1)$.
\end{enumerate}

\textcolor{darkaccent}{\textbf{Пример с доски.}} Игорь положил $750$ тыс. руб., через год добавил $450$ тыс. руб. Через два года сумма составила $1487{,}7$ тыс. руб. Найдём ставку:
\[
750k^2+450k=1487{,}7.
\]
\[
\begin{aligned}
750k^2+450k-1487{,}7&=0,\\
7500k^2+4500k-14877&=0,\\
D&=4500^2+4\cdot7500\cdot14877\\
 &=466\,560\,000=46\,656\cdot10\,000,\\
\sqrt D&=216\cdot100=21\,600,\\
k&=\frac{-4500+21600}{15000}=1{,}14.
\end{aligned}
\]
Следовательно, $r=100(1{,}14-1)=\boxed{14\%}$. Второй корень отрицателен и не соответствует росту вклада.

\sepheading{Как извлекать корни из больших дискриминантов}
\textbf{Приём 1. Разложение на множители.} На доске разобрали число $46\,656$:
\[
46\,656=2^6\cdot3^6,\qquad
\sqrt{46\,656}=2^3\cdot3^3=8\cdot27=216.
\]
Для произведения неотрицательных чисел используем $\sqrt{ab}=\sqrt a\sqrt b$. Поэтому
\[
\sqrt{466\,560\,000}=\sqrt{46\,656}\sqrt{10\,000}=216\cdot100=21\,600.
\]

\textbf{Приём 2. Оценка между соседними десятками.} Если требуется вычислить $\sqrt{1849}$, то
\[
40^2=1600<1849<2500=50^2,
\]
значит, искомый целый корень находится между $40$ и $50$. Поскольку число оканчивается на $9$, последняя цифра его целого корня может быть $3$ или $7$. Проверяем $43^2=1849$, значит,
\[
\boxed{\sqrt{1849}=43}.
\]
\textbf{Проверка из занятия:} $\sqrt{5184}=72$, поскольку $72^2=5184$.

\textbf{Пример с большим дискриминантом.} Для вклада $850$ тыс. руб., пополнения $350$ тыс. руб. и итоговой суммы $1549{,}76$ тыс. руб.:
\[
\begin{aligned}
850k^2+350k-1549{,}76&=0,\\
D=350^2+4\cdot850\cdot1549{,}76&=5\,391\,684,\\
\sqrt D=\sqrt{(2\cdot3^3\cdot43)^2}&=2322,\\
k=\frac{-350+2322}{1700}&=1{,}16.
\end{aligned}
\]
\key{Ответ:} $16\%$ годовых.

\clearpage
\section*{Тренировочный блок --- Путь А}
\noindent\textcolor{accent}{\rule{\linewidth}{0.45pt}}
\textit{10 заданий: от распознавания схем до восстановления процентной ставки. Решения оформляйте последовательно.}

\begin{enumerate}[label=\textbf{\arabic*.}]
\item \textbf{Определите схему кредитования} в каждом случае: \\
а) ежегодно тело долга уменьшается на $240$ тыс. руб.; \\
б) ежегодно заёмщик вносит одну и ту же сумму; \\
в) каждый год остаток долга становится на $20\%$ меньше предыдущего (о равных платежах ничего не сказано).

\item Кредит $1\,200\,000$ руб. выдан на $4$ года под $10\%$ годовых. Тело долга погашается равными годовыми частями. Найдите \textbf{общую сумму выплат}.

\item Кредит $900\,000$ руб. выдан на $3$ года под $12\%$ годовых по дифференцированной схеме. Найдите \textbf{первый и последний платежи}.

\item Кредит $600\,000$ руб. выдан на $5$ лет под $8\%$ годовых с равномерным погашением основного долга. Найдите \textbf{переплату в рублях и в процентах от суммы кредита}.

\item Вклад открыт на два года под $10\%$ годовых. Первоначально внесли $500$ тыс. руб., после первого года пополнили его на $100$ тыс. руб. Найдите сумму \textbf{в конце второго года}.

\item На вклад под $5\%$ годовых внесли $400$ тыс. руб. Через год, после начисления процентов, сняли $60$ тыс. руб. Сколько денег останется \textbf{через два года}?

\item \textbf{Вычислите без калькулятора:} $\sqrt{5184}$, $\sqrt{1849}$ и $\sqrt{46656}$. Кратко поясните способ для каждого числа.

\item Первоначально на вклад внесли $300$ тыс. руб., через год пополнили его на $70$ тыс. руб. В конце второго года получилось $516$ тыс. руб. Найдите \textbf{годовую процентную ставку}.

\item Вклад на два года: первоначальная сумма $700$ тыс. руб., пополнение через год $300$ тыс. руб., итог $1251{,}72$ тыс. руб. Определите \textbf{годовую процентную ставку}, составив квадратное уравнение.

\item Дифференцированный кредит на $4$ года составил $1\,800\,000$ руб. Всего банку выплатили $2\,250\,000$ руб. Найдите \textbf{годовую процентную ставку, первый и последний платежи}.
\end{enumerate}

\clearpage
\section*{Ответы к тренировочному блоку}
\noindent\textcolor{accent}{\rule{\linewidth}{0.45pt}}
\vspace{4mm}
\renewcommand{\arraystretch}{1.35}
\begin{tabularx}{0.99\linewidth}{@{}p{11mm}X@{}}
\toprule
\textbf{№} & \textbf{Ответ}\\
\midrule
1 & а) дифференцированная; б) аннуитетная; в) произвольная.\\
2 & $1\,500\,000$ руб.\\
3 & Первый --- $408\,000$ руб., последний --- $336\,000$ руб.\\
4 & $144\,000$ руб.; $24\%$.\\
5 & $715$ тыс. руб.\\
6 & $378$ тыс. руб.\\
7 & $72$; $43$; $216$.\\
8 & $20\%$ годовых.\\
9 & $14\%$ годовых.\\
10 & $10\%$ годовых; первый платёж $630\,000$ руб.; последний --- $495\,000$ руб.\\
\bottomrule
\end{tabularx}

\vspace{13mm}
\begin{center}
{\large\color{darkaccent}Краткая памятка}\par
\vspace{5mm}
\begin{tabularx}{0.95\linewidth}{@{}p{0.41\linewidth}X@{}}
\toprule
\textbf{Увидели в задаче} & \textbf{Что делать}\\
\midrule
Равные платежи & Проверить аннуитетную схему.\\
Равное снижение долга в рублях & Применить формулы дифференцированного кредита.\\
Снижение долга на одинаковый процент & Рассмотреть произвольную схему.\\
Двухлетний вклад с пополнением & $Sk^2+Ak=F$, где $k=1+r/100$.\\
Большой дискриминант & Вынести квадратный множитель из-под корня или разложить на простые множители.\\
\bottomrule
\end{tabularx}
\end{center}

\vspace{4mm}
\textcolor{darkaccent}{\textbf{Самопроверка перед сдачей работы}}
\begin{itemize}[itemsep=0pt,topsep=2pt]
\item Различены равные платежи и равные части основного долга.
\item Проценты кредита начислены на остаток; учтён момент взноса на вклад.
\item Использованы $p=r/100$ и $k=1+r/100$.
\item Десятичные дроби устранены умножением всех членов уравнения.
\item Проверены дискриминант, допустимость корня и единицы измерения.
\end{itemize}

\vfill
\begin{center}
{\small\color{softgray}Лёвин Артём Александрович эксклюзивно для Анастасии Павловой}
\end{center}
\end{document}
